\documentclass[12pt]{article}

\usepackage[portuguese]{babel}
\usepackage{float}
\usepackage{listings}
\usepackage{indentfirst}
\usepackage{amsfonts, amssymb, amsmath}
\usepackage{graphicx}
\usepackage{enumitem}
\usepackage[colorlinks=true, allcolors=blue]{hyperref}
\usepackage[letterpaper,
            top=1.5cm,
            bottom=1.5cm,
            left=2.75cm,
            right=2.75cm,
            marginparwidth=1.75cm]{geometry}

\title{Lista 1 -- Conceitos Básicos de LFA}
\author{Guilherme Lopes}
\date{\today}

\begin{document}
\maketitle

    % \begin{itemize}[label=R:]
    %     \item 
    % \end{itemize}
    
\begin{enumerate}
    \item O que é alfabeto?
    \begin{itemize}[label=R:]
        \item Um alfabeto é um conjunto de simbolos finito e não vazio
    \end{itemize}
    \item Defina o conceito de cadeia.
    \begin{itemize}[label=R:]
        \item Cadeia ou palavra é uma sequência de simbolos presentes em um 
        alfabeto
    \end{itemize}
    \item Defina o conceito de linguagem e mostre um exemplo.
    \begin{itemize}[label=R:]
        \item Linguagem é um conjunto de cadeias presentes dentro de um 
        fechamento de alfabeto
    \end{itemize}
    \item O que é fechamento de um alfabeto?
    \begin{itemize}[label=R:]
        \item O fechamento de um alfabeto é a união de todas as infinitas 
        exponenciações desse alfabeto (a partir de $\Sigma^0$). 
    \end{itemize}
    \item Como se pode descrever uma linguagem formal?
    \begin{itemize}[label=R:]
        \item É um subconjunto de cadeias presente dentro do fechamento de um 
        alfabeto.

    \end{itemize}
    Resposta do professor: Uma linguagem formal pode ser descrita 
    por um modelo reconhecedor ou por um modelo gerador. Um modelo reconhecedor 
    é um modelo matemático capaz de percorrer (“varrer”) uma cadeia de símbolos 
    construída a partir de um alfabeto e, ao final desta “varredura”, identificar 
    se esta cadeia faz parte ou não da linguagem descrita por ele. Neste 
    contexto, a linguagem descrita pelo modelo corresponde ao conjunto formado 
    por todas as cadeias que ele aceita. Já um modelo gerador é um modelo capaz 
    de gerar (produzir) as cadeias que fazem parte de uma linguagem definida a 
    partir de um alfabeto. Neste contexto, a linguagem descrita pelo modelo 
    corresponde ao conjunto de todas as cadeias que ele é capaz de gerar 
    (produzir). 
    \item Fale sobre aplicações de LFA.
    \begin{itemize}[label=R:]
        \item LFA pode ser aplicada na \dots
    \end{itemize}
    Resposta do professor: É difícil descrever todo o universo de aplicações 
    nas quais se pode utilizar os modelos estudados em LFA. Entretanto, entre 
    as principais aplicações, pode-se destacar:
    \begin{itemize}
        \item análise de linguagens de programação
        \begin{itemize}
            \item léxica;
            \item sintática;
        \end{itemize}
        \item modelos de sistemas biológicos;
        \item desenho de hardware;
        \item relacionamentos com linguagens naturais;
        \item processamento de imagens ou visão computacional (reconhecimento de \\
    padrões);
    \end{itemize}

    \item Defina o conceito de subpalavra.
    \begin{itemize}[label=R:]
        \item Subpalavra é uma cadeia que está no meio de outras duas, um 
        exemplo, é a palavra $x$, que é uma subpalavra em $txw$. 
    \end{itemize}
    \item Prove que se uma cadeia $x$ é prefixo de uma cadeia $y$ e $y$ também é 
    prefixo de $x$, então $x$ e $y$ são iguais.
    \begin{itemize}[label=R:]
        \item Vamos considerar as palavras $x=yz$ e $y=xu$, com $x,y,z, u \in 
        \Sigma^*$.

        Por substituição, podemos dizer que $x=xuz$. 

        Como $x=xuz$, consideramos $uz=\lambda$, assim, $x=x\lambda=x$.

        Se $uz=\lambda$, $u=z=\lambda$, e $y=xu$, logo $y=x\lambda=x$.

        Assim, para $x$ ser prefixo de $y$ e $y$ ser prefixo de $x$, $y=x$.

    \end{itemize}
    \item Prove que se uma cadeia $x$ é prefixo de uma cadeia $y$ e $y$ é prefixo 
    de uma cadeia $z$, então $x$ é prefixo de $z$.
    \begin{itemize}[label=R:]
        \item Considerando $y=xa$, $z=yb$, com $x,y,z,a,b \in \Sigma^*$.
        
        Substituindo o $y$ na palavra $z$, obtemos $z=xab$.

        Como $x$ também é uma cadeia, podemos dizer que $x$ também é prefixo de 
        $z$.
    \end{itemize}
    \item Dados $L_1={a, ab}$ e $L_2={\lambda, a, ba}$, linguagens sobre 
    ${a, b}$, determine:
    \begin{enumerate}[label=\alph*.]
        \item $L_1 \cup L_2$
        \begin{itemize}[label=R:]
            \item $\{\lambda,a,ab,ba\}$
        \end{itemize}
        \item $L_1 \cap L_2$
        \begin{itemize}[label=R:]
            \item $\{a\}$
        \end{itemize}
        \item $L_1 - L_2$
        \begin{itemize}[label=R:]
            \item $\{ab\}$
        \end{itemize}
        \item $L_2 - L_1$
        \begin{itemize}[label=R:]
            \item $\{\lambda,ba\}$
        \end{itemize}
        \item $L_1.L_2$
        \begin{itemize}[label=R:]
            \item $\{a,aa,aba,ab,aba,abba\}$
        \end{itemize}
        \item $L_2.L_1$
        \begin{itemize}[label=R:]
            \item $\{a,ab,aa,aab,baa,baab\}$
        \end{itemize}
        \item $(L_1)^2 = L_1.L_1$
        \begin{itemize}[label=R:]
            \item $\{aa,aab,aba,abab\}$
        \end{itemize}
        \item $(L_2)^2 = L_2.L_2$
        \begin{itemize}[label=R:]
            \item $\{\lambda,a,ba,aa,aba,baa,baba\}$
        \end{itemize}
        \item $\bar{L_1}$
        \begin{itemize}[label=R:]
            \item $\bar{L_1}=\Sigma^*-L_1={a,b}^*-{a,ab}$
        \end{itemize}
    \end{enumerate}
\end{enumerate}

\end{document}
